% ============================================================
% 分片：s02a-fieldtheory  对应原文 PDF 第 4--9 页（印刷页码 4--9）
% 负责章节：2. FIELD THEORY BACKGROUND（前半，§2 标题起至第 9 页末）
% 译者：t02-s2a   审校状态：初稿
% 审校：t05-s32 已完成逐式核对（修正 0 处）
% ============================================================
% [说明: 第 4 页自 §2 标题开始；该页 §2 标题之前的引言尾段属 s01-intro.tex]
% [脚注说明: 本分片脚注由 LaTeX 自动编号，与原文第 7 页的星号脚注相比编号整体偏移，属正常]

% ==== 原文 PDF 第 4 页（印刷页 4）====

\section{场论背景}

\begin{keytakeaway}
\textbf{核心物理图像：} 二维 $N=2$ 超对称理论可由四维 $N=1$ 理论经维数约化得到，因而其超场代数、$\sigma$ 矩阵与费米子约定都继承自四维；超对称的表示与规范不变性完全编码在超空间微分算符的代数中。Wess--Zumino 规范把矢量超场的多余分量规范掉，只留下规范场、规范微子（gaugino）与辅助场 $D$。二维区别于四维的新结构是：除手征超场外还存在扭手征超场（twisted chiral），而规范场的场强本身就是一个扭手征超场。
\par\textbf{数学工具：} 超空间坐标 $x^m,\theta^\alpha,\bar\theta^{\dot\alpha}$ 上的超对称生成元 $Q_\alpha,\bar Q_{\dot\alpha}$ 与协变导数 $D_\alpha,\bar D_{\dot\alpha}$（及其规范版本 $\mathcal D_\alpha,\bar{\mathcal D}_{\dot\alpha},\mathcal D_m$）；手征与反手征超场的 $\theta$ 展开 (2.4)--(2.5)；矢量超场的 Wess--Zumino 规范形 (2.11)；维数约化后的分量变换律 (2.12)、(2.14)；以及扭手征场强 $\Sigma=\frac{1}{2\sqrt2}\{\bar{\mathcal D}_+,\mathcal D_-\}$ 及其分量展开 (2.16)。
\par\textbf{主要结论：} 本半节固定后续分析所需的全部约定与带号公式：$\epsilon^{12}=-\epsilon_{12}=1$ 与 $\sigma^m_{\alpha\dot\alpha}$ 的表示取自 [10]；Wess--Zumino 规范下矢量多重态的分量为 $v_m,\lambda_\alpha,\bar\lambda_{\dot\alpha},D$；二维的基本规范不变场强是扭手征超场 $\Sigma$。原文明确指出，这些详细公式对读懂本文并非必需，列出它们是为了备查并使叙述自足。
\end{keytakeaway}

本节给出后续分析所需的场论背景。本节大多数详细公式在阅读本文时并非必需，列出它们是为了便于查阅，并使叙述尽量自足。

二维 $N=2$ 超对称可由四维时空的 $N=1$ 超对称经维数约化（dimensional reduction）得到。这方面已有若干不错的教科书 [10--14]。作者对四维超场的约定采用 [10] 中的约定。

在以 $x^m,\theta^\alpha,\bar\theta^{\dot\alpha}$ 为坐标的超空间中，超对称由下列算符几何地实现
% [跨页接续: 原文第4页/第5页；本句的算符展示式 (2.1) 排在第 5 页]

% ==== 原文 PDF 第 5 页（印刷页 5）====
\label{op:5}

\begin{equation}
\begin{aligned}
Q_\alpha &= \frac{\partial}{\partial\theta^\alpha} - i\sigma^m_{\alpha\dot\alpha}\bar\theta^{\dot\alpha}\frac{\partial}{\partial x^m}\\
\bar Q_{\dot\alpha} &= -\frac{\partial}{\partial\bar\theta^{\dot\alpha}} + i\sigma^m_{\alpha\dot\alpha}\theta^\alpha\frac{\partial}{\partial x^m}
\end{aligned}
\tag{2.1}
\end{equation}

$\alpha$ 与 $\dot\alpha$ 是自旋指标的两种手征性；此外还记 $\psi_\alpha=\epsilon_{\alpha\beta}\psi^\beta$、$\psi^\alpha=\epsilon^{\alpha\beta}\psi_\beta$，其中 $\epsilon$ 为反对称张量，$\epsilon^{12}=-\epsilon_{12}=1$；点指标情形相同。张量 $\sigma^m_{\alpha\dot\alpha}$ 采用 [10] 中所用的表示

% [Tutor Note: 指标升降的符号最易错：由 $\epsilon^{12}=-\epsilon_{12}=1$ 得 $\psi^1=\epsilon^{12}\psi_2=\psi_2$、$\psi^2=\epsilon^{21}\psi_1=-\psi_1$。第 8 页把分量记为 $(\psi^1,\psi^2)=(\psi^-,\psi^+)$、$(\psi_1,\psi_2)=(\psi_-,\psi_+)$ 并注明 $\psi^-=\psi_+$、$\psi^+=-\psi_-$，正是这条约定的推论；核对 (2.4) 中 $\theta^\alpha\psi_\alpha$ 的缩并以及 (2.12)、(2.14) 中 $\epsilon_\pm$ 与 $\psi_\pm$ 的搭配时，不可把上下标随意互换。]

\begin{equation}
\begin{aligned}
\sigma^0 &= \begin{pmatrix} -1 & 0 \\ 0 & -1 \end{pmatrix}\\
\sigma^1 &= \begin{pmatrix} 0 & 1 \\ 1 & 0 \end{pmatrix}\\
\sigma^2 &= \begin{pmatrix} 0 & -i \\ i & 0 \end{pmatrix}\\
\sigma^3 &= \begin{pmatrix} 1 & 0 \\ 0 & -1 \end{pmatrix}
\end{aligned}
\tag{2.2}
\end{equation}

% [Tutor Note: 注意 $\sigma^0=-\mathbb{1}$（而非 $+\mathbb{1}$），且 $\sigma^0,\sigma^3$ 均为对角阵；配合度规 $\eta^{mn}=\mathrm{diag}(-1,1,1,1)$ 的 $(-,+,+,+)$ 号差，这套约定同时固定了 $\sigma^m_{\alpha\dot\alpha}$ 的上下标位置。若混用取 $\sigma^0=+\mathbb{1}$ 的教材约定，(2.12) 中 $\partial_0\pm\partial_1$ 与 $\epsilon_\pm$ 的搭配以及 $v_{01}$ 项的符号都会整体改变。]

度规为 $\eta^{mn}=\mathrm{diag}(-1,1,1,1)$。式 (2.1) 的超对称生成元与下列算符对易

% [Tutor Note: 比较 (2.1) 与 (2.3) 可得 $Q_\alpha=D_\alpha-2i\sigma^m_{\alpha\dot\alpha}\bar\theta^{\dot\alpha}\partial_m$、$\bar Q_{\dot\alpha}=\bar D_{\dot\alpha}+2i\sigma^m_{\alpha\dot\alpha}\theta^\alpha\partial_m$：超对称生成元与超协变导数只相差一个正比于 $\partial_m$ 的项。原文说生成元与 $D$ 算符“对易”，即指二者在对易关系中的这种相容性——这正是手征条件 $\bar D_{\dot\alpha}\Phi=0$ 与超对称变换可以同时定义的前提。]

\begin{equation}
\begin{aligned}
D_\alpha &= \frac{\partial}{\partial\theta^\alpha} + i\sigma^m_{\alpha\dot\alpha}\bar\theta^{\dot\alpha}\frac{\partial}{\partial x^m}\\
\bar D_{\dot\alpha} &= -\frac{\partial}{\partial\bar\theta^{\dot\alpha}} - i\sigma^m_{\alpha\dot\alpha}\theta^\alpha\frac{\partial}{\partial x^m}
\end{aligned}
\tag{2.3}
\end{equation}

这些算符用于书写拉格朗日量。

最简单的超场类型是手征超场（chiral superfield）$\Phi$，它满足 $\bar D_{\dot\alpha}\Phi=0$，并可展开为

\begin{equation}
\Phi(x,\theta)=\phi(y)+\sqrt2\theta^\alpha\psi_\alpha(y)+\theta^\alpha\theta_\alpha F(y),
\tag{2.4}
\end{equation}

其中 $y^m=x^m+i\theta^\alpha\sigma^m_{\alpha\dot\alpha}\bar\theta^{\dot\alpha}$。$\Phi$ 的复共轭是反手征多重态（antichiral multiplet）$\bar\Phi$，满足 $D_\alpha\bar\Phi=0$，其展开为

% ==== 原文 PDF 第 6 页（印刷页 6）====
\label{op:6}

\begin{equation}
\bar\Phi=\bar\phi(\bar y)+\sqrt2\bar\theta_{\dot\alpha}\bar\psi^{\dot\alpha}(\bar y)+\bar\theta_{\dot\alpha}\bar\theta^{\dot\alpha}F(\bar y),
\tag{2.5}
\end{equation}
% [原文疑误: (2.5) 既是 (2.4) 的复共轭，末项应为 $\bar\theta_{\dot\alpha}\bar\theta^{\dot\alpha}\bar F(\bar y)$；原文印刷为不含横线的 $F$，而同式中 $\bar\phi$、$\bar\psi$ 均带横线，故判为漏排。建议补为 $\bar F$（本分片照抄原文 $F$）]

其中 $\bar y^m=x^m-i\theta^\alpha\sigma^m_{\alpha\dot\alpha}\bar\theta^{\dot\alpha}$。一般地，以大写字母 $\Phi,S,P,\dots$ 标记的手征场，其 $\theta=0$ 分量将用相应的小写字母 $\phi,s,p,\dots$ 表示。

为构造规范理论，需要在超空间中引入规范场，把微分算符 $D_\alpha$、$\bar D_{\dot\alpha}$ 与 $\partial_m=\partial/\partial x^m$ 替换为规范协变导数 $\mathcal D_\alpha$、$\bar{\mathcal D}_{\dot\alpha}$ 与 $\mathcal D_m$。然而必须对超空间规范场施加很强的限制。为使带电荷的手征超场得以存在，需要方程 $\bar{\mathcal D}_{\dot\alpha}\Phi=0$ 可积。因此要求

% [Tutor Note: (2.6) 的必要性：若 $\{\bar{\mathcal D}_{\dot\alpha},\bar{\mathcal D}_{\dot\beta}\}\neq0$，则对 $\bar{\mathcal D}_{\dot\alpha}\Phi=0$ 再作用一次 $\bar{\mathcal D}_{\dot\beta}$ 会给出含 $\mathcal D_m\Phi$ 的非零项，手征条件自相矛盾。进一步的约束 (2.8) 把其余反对易子固定为 $-2i\sigma^m_{\alpha\dot\alpha}\mathcal D_m$，其含义是超空间规范场没有独立于 $V$ 的额外自由度；约化到分量即为 (2.11) 的展开。]

\begin{equation}
0=\{\bar{\mathcal D}_{\dot\alpha},\bar{\mathcal D}_{\dot\beta}\}=\{\mathcal D_\alpha,\mathcal D_\beta\}.
\tag{2.6}
\end{equation}

这些条件保证，在适当的部分规范固定下，

\begin{equation}
\begin{aligned}
\mathcal D_\alpha &= e^{-V}D_\alpha e^{V}\\
\bar{\mathcal D}_{\dot\alpha} &= e^{V}\bar D_{\dot\alpha}e^{-V},
\end{aligned}
\tag{2.7}
\end{equation}

其中 $V$ 是超空间上的实李代数值函数；$V$ 称为矢量超场（vector superfield）。此外还施加进一步约束

\begin{equation}
\{\mathcal D_\alpha,\bar{\mathcal D}_{\dot\alpha}\}=-2i\sigma^m_{\alpha\dot\alpha}\mathcal D_m,
\tag{2.8}
\end{equation}

使超空间规范场完全由 $V$ 确定。

仍存在一个剩余规范不变性；对规范群 $U(1)$——此时 $V$ 只是超空间上的单个实函数——其形式为

\begin{equation}
V\to V+i(\Lambda-\bar\Lambda),
\tag{2.9}
\end{equation}

其中 $\Lambda$ 为手征超场。电荷为 $Q$ 的手征超场 $\Phi$ 在该剩余规范不变性下按如下方式变换

% ==== 原文 PDF 第 7 页（印刷页 7）====
\label{op:7}

\begin{equation}
\Phi\to\exp(-iQ\Lambda)\cdot\Phi.
\tag{2.10}
\end{equation}

可以通过部分固定这一剩余规范不变性（即进入所谓 Wess--Zumino 规范）把 $V$ 写成如下形式

\begin{equation}
V=-\theta^\alpha\sigma^m_{\alpha\dot\alpha}\bar\theta^{\dot\alpha}v_m+i\theta^\alpha\theta_\alpha\bar\theta_{\dot\alpha}\bar\lambda^{\dot\alpha}-i\bar\theta_{\dot\alpha}\bar\theta^{\dot\alpha}\theta^\alpha\lambda_\alpha+\frac{1}{2}\theta^\alpha\theta_\alpha\bar\theta_{\dot\alpha}\bar\theta^{\dot\alpha}D.
\tag{2.11}
\end{equation}

% [Tutor Note: Wess--Zumino 规范的消去机制：一般实矢量超场还含 $C$（与 $\theta$ 无关的分量）、$\chi_\alpha,\bar\chi_{\dot\alpha}$（$\theta$ 与 $\bar\theta$ 的一次项）以及 $M,N$（$\theta\theta$ 与 $\bar\theta\bar\theta$ 项）；利用手征规范参数 $\Lambda$（$\Lambda-\bar\Lambda$ 提供足够的玻色与费米参数）可把 $C,\chi,\bar\chi,M,N$ 全部规范掉，只留下 $v_m,\lambda_\alpha,\bar\lambda_{\dot\alpha}$ 与 $D$。$D$ 无动能项、只以代数形式出现，故可积掉。后文加入 Fayet--Iliopoulos 项（式 (2.27) 中的 $-rD$）后，$D$ 的代数运动方程使 $D$ 的真空值被 FI 参数 $r$ 平移，代回拉格朗日量即得 D 项势能，这正是后文判定相结构的出发点。]

在 Wess--Zumino 规范下，通常的规范不变性仍然保留。它对应于 $\Lambda=-a(x)$（$a$ 为只依赖 $x$ 的实函数），以及通常的变换律 $\Phi\to\exp(iQa)\Phi$、$v\to v-da$。

在四维中，阿贝尔（abelian）规范场的基本规范不变场强是 $[\mathcal D_\alpha,\mathcal D_m]$。在二维中还存在一个更基本的、稍后引入的不变量。

在 Wess--Zumino 规范下，理论的物理内容虽然相对清晰，但超对称变换律却相对复杂，因为超对称变换（由 $Q_\alpha$ 与 $\bar Q_{\dot\alpha}$ 生成）必须伴随规范变换才能保持 Wess--Zumino 规范。当维数约化到二维时，变换律变得更加复杂。\footnote{全部约定均取自 [10] 的四维约定经维数约化所得。作者采用这一途径，是因为将要研究的模型由四维约化而来十分自然。此外，关于二维矢量超场的分量展开似乎并不存在任何标准约定；二维的手征超场常用与本文相差因子 $\sqrt2$ 的约定来描述，这些因子可以通过对场作相当显然的重新标度而吸收。} 下面把约化后的 Wess--Zumino 规范变换律写出，以供参考；但作者愿立即向读者保证：不熟悉下列公式并不妨碍读懂本文。

作约化时，取各场与 $x^1$、$x^2$ 无关。于是规范场在 $x^1$、$x^2$ 方向的分量 $v_1$ 与 $v_2$，连同所有其他场，都只是 $x^0$、$x^3$ 的函数；记 $\sigma=(v_1-iv_2)/\sqrt2$、$\bar\sigma=(v_1+iv_2)/\sqrt2$。
% [跨页接续: 原文第7页/第8页；该句自第 7 页末行续至第 8 页]

% ==== 原文 PDF 第 8 页（印刷页 8）====
\label{op:8}

对余下的时空坐标采用二维记号，记 $x^0=y^0$、$x^3=y^1$。维数约化之后，宜将费米子分量标记为 $(\psi^1,\psi^2)=(\psi^-,\psi^+)$ 与 $(\psi_1,\psi_2)=(\psi_-,\psi_+)$（于是 $\psi^-=\psi_+$、$\psi^+=-\psi_-$），点分量同此。矢量多重态在参数为 $\epsilon_\pm,\bar\epsilon_\pm$ 的超对称变换下的维数约化变换律为（由 [10] 第 50 页的公式约化而得）

\begin{equation}
\begin{aligned}
\delta v_m &= i\bar\epsilon\sigma_m\lambda+i\epsilon\sigma_m\bar\lambda\\
\delta\sigma &= -i\sqrt2\bar\epsilon_+\lambda_--i\sqrt2\epsilon_-\bar\lambda_+\\
\delta\bar\sigma &= -i\sqrt2\epsilon_+\bar\lambda_--i\sqrt2\bar\epsilon_-\lambda_+\\
\delta D &= -\bar\epsilon_+(\partial_0-\partial_1)\lambda_+-\bar\epsilon_-(\partial_0+\partial_1)\lambda_-+\epsilon_+(\partial_0-\partial_1)\bar\lambda_++\epsilon_-(\partial_0+\partial_1)\bar\lambda_-\\
\delta\lambda_+ &= i\epsilon_+D+\sqrt2(\partial_0+\partial_1)\bar\sigma\epsilon_- -v_{01}\epsilon_+\\
\delta\lambda_- &= i\epsilon_-D+\sqrt2(\partial_0-\partial_1)\sigma\epsilon_+ +v_{01}\epsilon_-\\
\delta\bar\lambda_+ &= -i\bar\epsilon_+D+\sqrt2(\partial_0+\partial_1)\sigma\bar\epsilon_- -v_{01}\bar\epsilon_+\\
\delta\bar\lambda_- &= -i\bar\epsilon_-D+\sqrt2(\partial_0-\partial_1)\bar\sigma\bar\epsilon_+ +v_{01}\bar\epsilon_-,
\end{aligned}
\tag{2.12}
\end{equation}

% [Tutor Note: (2.12) 中的 $\pm$ 是二维轻锥标签，不是超空间自旋指标；它与第 8 页 $(\psi^1,\psi^2)=(\psi^-,\psi^+)$、$(\psi_1,\psi_2)=(\psi_-,\psi_+)$ 的记法并用，故出现 $\psi^-=\psi_+$、$\psi^+=-\psi_-$ 这类“换指标即变号”的关系。另外 $\sigma=(v_1-iv_2)/\sqrt2$ 与 $\bar\sigma$ 是互为复共轭的复标量（不是 $\sigma^m$ 矩阵）。逐项核对时须看三对共轭量是否配对正确：$\sigma$ 与 $\bar\sigma$、$\epsilon$ 与 $\bar\epsilon$、$\lambda$ 与 $\bar\lambda$；任一处错配都会使 $v_{01}$ 项的符号与 (2.12) 不符。]

其中 $v_{01}=\partial_0v_1-\partial_1v_0$。

存在规范场时，处于规范群某个给定表示中的手征超场 $\Phi$ 是满足 $\bar{\mathcal D}_{\dot\alpha}\Phi=0$ 的超场，其中 $\bar{\mathcal D}_{\dot\alpha}$ 里的协变导数取相应表示。若记

\begin{equation}
\Phi=e^{V}\Phi_0,
\tag{2.13}
\end{equation}

则 $\Phi_0$ 满足 $\bar{\mathcal D}_{\dot\alpha}\Phi_0=0$，并具有 (2.4) 形式的 $\theta$ 展开。

% ==== 原文 PDF 第 9 页（印刷页 9）====
\label{op:9}

这些分量场的变换律为

\begin{equation}
\begin{aligned}
\delta\phi &= \sqrt2(\epsilon_+\psi_--\epsilon_-\psi_+)\\
\delta\psi_+ &= i\sqrt2(D_0+D_1)\phi\bar\epsilon_-+\sqrt2\epsilon_+F-2Q\phi\bar\sigma\bar\epsilon_+\\
\delta\psi_- &= -i\sqrt2(D_0-D_1)\phi\bar\epsilon_++\sqrt2\epsilon_-F+2Q\phi\sigma\bar\epsilon_-\\
\delta F &= -i\sqrt2\bar\epsilon_+(D_0-D_1)\psi_+-i\sqrt2\bar\epsilon_-(D_0+D_1)\psi_-\\
&\quad +2Q(\bar\epsilon_+\bar\sigma\psi_-+\bar\epsilon_-\sigma\psi_+)+2iQ\phi(\bar\epsilon_-\bar\lambda_+-\bar\epsilon_+\bar\lambda_-).
\end{aligned}
\tag{2.14}
\end{equation}

% [Tutor Note: 辅助场的判据：$\delta F$ 与 (2.12) 中的 $\delta D$ 都不含 $\partial F$、$\partial D$ 本身，只含其他分量及其导数，因此 $F$、$D$ 在拉格朗日量中只作代数出现，运动方程可直接解出并代回。再注意符号冲突：(2.14) 中的 $D_0\pm D_1$ 是含 $v_m$ 与电荷 $Q$ 的规范协变导数，而 $D_\alpha,\bar D_{\dot\alpha}$ 是超空间导数——原文以花体 $\mathcal D$ 区分规范协变导数，但 $D_0,D_1$ 仍用斜体，翻译时保留原符号，阅读时须靠上下文分辨。]

反手征多重态 $\bar\Phi$ 的变换律是上述各式的复共轭。

\subsection*{扭手征超场；规范场强}

相对于四维，二维中出现的一个新特点是：除满足 $\bar D_+\Phi=\bar D_-\Phi=0$ 的手征超场外，还可以有满足 $\bar D_+\Sigma=D_-\Sigma=0$ 的扭手征超场（twisted chiral superfield）[15--19]。同时含手征超场与扭手征超场的 $\sigma$ 模型颇为优美。由于镜像对称把手征多重态变为扭手征多重态，考虑同时含这两类多重态的适当模型很可能有助于理解镜像对称。

然而，对当前目的更为重要的是如下事实（[20] 在把 $N=2$ 对偶解释为阿贝尔镜像对称时用到了它）：在二维中，超空间规范场的基本规范不变场强是一个扭手征超场。该量为

\begin{equation}
\Sigma=\frac{1}{2\sqrt2}\{\bar{\mathcal D}_+,\mathcal D_-\}.
\tag{2.15}
\end{equation}

% [Tutor Note: $\Sigma$ 是四维场强 $W_\alpha=\bar D\bar D D_\alpha V$ 在二维的对应物，但带两个相反手征性的指标 $\bar{\mathcal D}_+\,\mathcal D_-$，这正是它能成为扭手征超场的原因：扭手征条件 $\bar{\mathcal D}_+\Sigma=\mathcal D_-\Sigma=0$ 与手征条件 $\bar{\mathcal D}_\pm\Phi=0$ 来自不同的子代数。由 (2.16) 可见 $\Sigma$ 的最低分量为复标量 $\sigma$，而 $\theta^+\theta^-$ 的系数是 $\sqrt2(D-iv_{01})$：二维中电磁场只有一个独立分量 $v_{01}$，它与辅助场 $D$ 复合成一个复辅助场，这是二维规范理论比四维更“刚性”的根源。]

按 Bianchi 恒等式，它被 $\bar{\mathcal D}_+$ 与 $\mathcal D_-$ 湮灭。在阿贝尔情形，有

\begin{equation}
\begin{aligned}
\Sigma &= \frac{1}{\sqrt2}\bar{\mathcal D}_+\mathcal D_-V=\sigma-i\sqrt2\theta^+\bar\lambda_+-i\sqrt2\bar\theta^-\lambda_-+\sqrt2\theta^+\bar\theta^-(D-iv_{01})\\
&\quad -i\bar\theta^-\theta^-(\partial_0-\partial_1)\sigma-i\theta^+\bar\theta^+(\partial_0+\partial_1)\sigma+\sqrt2\bar\theta^-\theta^+\theta^-(\partial_0-\partial_1)\bar\lambda_+\\
&\quad +\sqrt2\theta^+\bar\theta^-\bar\theta^+(\partial_0+\partial_1)\lambda_--\theta^+\bar\theta^-\theta^-\bar\theta^+(\partial_0^2-\partial_1^2)\sigma.
\end{aligned}
\tag{2.16}
\end{equation}
